\documentclass[10pt,a4paper]{amsart}
\usepackage[margin=19mm]{geometry}
\usepackage{amsmath,amssymb,mathtools}
\usepackage[scr=rsfs,cal=euler]{mathalfa}
\usepackage{xfrac}
\usepackage[unicode,hidelinks,bookmarks=false]{hyperref}
\newcommand{\ie}{i.e.\ }
\newcommand{\cf}{cf.\ }
\newcommand{\resp}{respectively\ }
\newcommand{\Subsection}{\subsection}
\providecommand{\nobreakdash}{\nobreak\hbox{-}}
\setlength{\emergencystretch}{2em}
\multlinegap0pt
\usepackage{tikz}
\usetikzlibrary{cd}
\usepackage{amssymb}
\usepackage{mathtools}
\newtheorem{lemma}{Lemma}
\newtheorem{proposition}{Proposition}
\title{The singular Hitchin fibration, cameral data, and representation theory}
\author{Alexander Fr\"{u}h}
\date{Source: arXiv:2602.00274v2 (CC BY 4.0)}
\begin{document}
\maketitle

\noindent\emph{Source and licence.} The singular Hitchin fibration, cameral data, and representation theory. Version arXiv:2602.00274v2 (CC BY 4.0), distributed by arXiv under the Creative Commons Attribution 4.0 International licence, http://creativecommons.org/licenses/by/4.0/, as recorded in the arXiv licence metadata for this exact version and shown on the abstract page https://arxiv.org/abs/2602.00274v2. Full licence notice: \texttt{licenses/arxiv-2602.00274-CC-BY-4.0.txt}.\par\smallskip

\noindent\textit{Editorial scope.} Counted unit: the article's proposition CameralHomprp (source0.tex lines 1058--1060). The article's complete original proof of it is delivered with it (source0.tex lines 1062--1105, one \texttt{proof} environment of 44 lines). No mathematics is written, repaired or re-derived: the statement and the proof are the article's own lines, in the article's own notation, and no retained line is substituted or re-flowed.  Retained uncounted as the source-local context the proof needs: lines 820--824; lines 963--965; lines 973--982; lines 2510--2515; lines 2570--2572.  Nothing was omitted from any retained span, so the package carries no omission record.  No retained line contains a citation, so the wrapper carries no bibliography and no bibliography entry.  The retained lines contain the article's own diagram code, which the wrapper typesets with the standard tikz packages; no image file is delivered and the package declares no asset dependency.  Editorial formatting only, in addition: an AMS \texttt{amsart} preamble that restates the article's own declarations and macros used by the retained lines, namely the environments \texttt{lemma}, \texttt{proposition} on separate counters, together with the standard packages those lines need.  The retained environments print in this document as Lemma 1, Proposition 1 in order of appearance, whereas the article prints the same items with the numbering of its own full document.  The complete native source of the article as distributed in the arXiv e-print for this exact version is bundled unchanged as \texttt{sources/193709/source0.tex}.
\begin{equation}\label{StackyChevalleyMapeqn}
\begin{tikzcd}
\rho_S:S \arrow[r] & {[S/G]} \arrow[r, "\rho_{S,G}"] & \mathcal B,
\end{tikzcd}
\end{equation}

\begin{lemma}\label{WLInvariancelem}
    For any scheme $X$ and map $X \rightarrow \mathcal B$, the $W_L$-action on $\mathfrak{z}$ lifts to a $W_L$-action on $X \times_{\mathcal B} \mathfrak{z}$.
\end{lemma}

\begin{proposition}\label{StackyGSprp}
    The diagram
    \begin{equation}\label{StackyGSeqn}
        \begin{tikzcd}
         \hat{S}^{reg} \arrow[r, "\hat{\chi}_S"] \arrow[d, "\hat{p}"] & \mathfrak{z} \arrow[d, "p"] \\
            S \arrow[r, "\rho_S"]                                       & \mathcal B              
        \end{tikzcd}
    \end{equation}
    is commutative and Cartesian, where $\rho_S:S \rightarrow \mathcal B$ is the $S$-Chevalley map defined in (\ref{StackyChevalleyMapeqn}).
\end{proposition}

\begin{equation}\label{GSDiagrameqn}
        \begin{tikzcd}
         \hat{S}^{reg} \arrow[r, "\hat{\chi}_S"] \arrow[d, "\hat{p}"] & \mathfrak{z} \arrow[d] \\
            S \arrow[r, "\chi_S"]                                       & \mathfrak{c}_S.              
        \end{tikzcd}
    \end{equation}

\begin{proposition}\label{ParabolicCentraliserprp}
    For $x \in (\mathfrak{z} +\overline{\mathcal O} +\mathfrak{n})^{reg}$, $\mathcal I_{S,x}^{sm} = C_P(x)$, as a subgroup of $\mathcal I_x = C_G(x)$.
\end{proposition}

\begin{proposition}\label{CameralHomprp}
There is a homomorphism of group schemes $\kappa_S:\mathcal I^{sm}_S \rightarrow \mathcal \rho_S^*\hat{\mathcal J}_S$, equivariant with respect to the actions of $G$ and $\mathbb{G}_m$, such that for any semisimple element $x \in S$, $\kappa_{S,x}: \mathcal I^{sm}_{S,x} \rightarrow \hat{\mathcal J}_{S, \rho_S(x)}$ realises the abelianisation of the group $\mathcal I^{sm}_{S,x}$.
\end{proposition}

\begin{proof} 
Let $P$ be a parabolic subgroup of $G$ with Levi factor $L$, and let $\mathfrak{r}$ be the solvable radical of $\mathfrak{p} = Lie(P)$. Let $\hat{p}:G \times^P\mathfrak{r}^{reg} \rightarrow S$ be the generalised Grothendieck-Springer map defined in (\ref{GSDiagrameqn}); by Proposition \ref{StackyGSprp}, we have that $\rho_S^* \Pi_S = \hat{p}_*(\bar{Z} \times (G \times^P \mathfrak{r}^{reg}))$. To construct a homomorphism $\kappa_S:\mathcal I^{sm}_S \rightarrow \mathcal \rho_S^*\Pi_S$, it is equivalent by adjunction to construct a homomorphism
\begin{equation}\label{CameralHomLifteqn}
    \hat{\kappa}_S:\hat{p}^*\mathcal I^{sm}_S \rightarrow \bar{Z} \times (G \times^P \mathfrak{r}^{reg}).
\end{equation}

There is a group scheme $\mathcal P$ over $G \times^P \mathfrak{r}^{reg}$ pulled back from the universal parabolic subgroup scheme of the constant group $G$ over $G/P$; in particular, for a $\mathbb{C}$-point $P(g,x)$ of $G \times^P \mathfrak{r}^{reg}$, the fibre of $\mathcal P$ over this point is $I_g(P)$. By Proposition \ref{ParabolicCentraliserprp}, there is an inclusion $\mathcal I_S^{sm} \hookrightarrow \mathcal P$. Moreover, there is a homomorphism from $\mathcal P$ to the constant group $\bar{Z}$ over $G \times^P \mathfrak{r}^{reg}$, which over the $\mathbb{C}$-point $P(g,x)$ is given by the composition
\begin{equation}\label{CameralHom1eqn}
\begin{tikzcd}
gPg^{-1} \arrow[r, "I_{g^{-1}}"] & P \arrow[r] & P/P^{der} = \bar{Z}. 
\end{tikzcd}
\end{equation}
To see that this is well-defined, we observe that if $gP = g'P$, then $I_g$ and $I_{g'}$ differ by conjugation by an element of $p$, and so define the same map after projection to $\bar{Z}$. Hence, we have a well-defined homomorphism $\hat{\kappa}_S$, and thus a well-defined homomorphism $\kappa_S:\mathcal I^{sm}_S \rightarrow \rho_S^*\Pi_S$. It is clear from the construction that $\kappa_S$ is equivariant with respect to the actions of $G$ and $\mathbb{G}_m$.

Since $\mathcal I_S^{sm}$ is smooth, to prove that $\kappa_S$ takes its image in the closed subgroup $\rho_S^* \hat{\mathcal J}_S$ of $\rho_S^*\Pi_S$, it suffices to do so over the dense open subset of semisimple elements $S^{ss} \subseteq S$. If we let $\mathfrak{z}^{rs}$ be the open subset of $\mathfrak{z}$ defined by
\begin{equation}\label{WeylLeviActioneqn}
    \mathfrak{z}^{rs} = \{ x \in \mathfrak{g} \, | \,C_G(x) = L \},
\end{equation}
then $S^{ss} = Ad(G)(\mathfrak{z}^{rs})$. By $G$-equivariance, $\kappa_S|_{S^{ss}}$ is determined by its value on $\mathfrak{z}^{rs}$, so in particular, $\kappa_S|_{S^{ss}}$ takes its image in $\rho_S^* \hat{\mathcal J}_S|_{S^{ss}}$ if and only if $\kappa_S|_{\mathfrak{z}^{rs}}$ takes its image in $\rho_S^* \hat{\mathcal J}_S|_{\mathfrak{z}^{rs}}$. 

There is a Cartesian diagram
\begin{equation}\label{CameralHom2eqn}
\begin{tikzcd}
W_L \times \mathfrak{z}^{rs} \arrow[d, "\pi_2"] \arrow[r, "\hat{\iota}"] & G \times^P \mathfrak{r}^{reg} \arrow[d, "\hat{p}"] \\
\mathfrak{z}^{rs} \arrow[r, hook]                                 & S                                                   
\end{tikzcd}
\end{equation}
where $\hat{\iota}$ is defined on $\mathbb{C}$-points by
\begin{equation}\label{CameralHom3eqn}
    (nL,x) \mapsto P(n^{-1},Ad_n(x)), 
\end{equation}
for $nL \in N_G(L)/L = W_L$ and $x \in \mathfrak{z}^{rs}$. We can identify the pullback of $\hat{\kappa}_S$ under $\hat{\iota}$ with the map
\begin{equation}\label{CameralHom4eqn}
    L \times (W_L \times \mathfrak{z}^{rs}) \rightarrow \bar{Z} \times (W_L \times \mathfrak{z}^{rs}) 
\end{equation}
which over $\{w\} \times \mathfrak{z}^{rs}$ is given by
the constant group homomorphism
\begin{equation}\label{CameralHom5eqn}
\begin{tikzcd}
L \arrow[r] & L/L^{der} = \bar{Z} \arrow[r, "w"] & \bar{Z}.
\end{tikzcd}  
\end{equation}
This is $W_L$-equivariant, where the $W_L$-action on $L \times (W_L \times \mathfrak{z}^{rs}) \cong \mathcal I_S^{sm} \times_{\mathcal B} \mathfrak{z}$ is given by Lemma \ref{WLInvariancelem}; hence $\kappa_S|_{\mathfrak{z}^{rs}}$ takes its image in $\rho_S^* \hat{\mathcal J}_{\mathfrak{z}^{rs}}$. This also gives the required characterisation of $\kappa_{S,x}$ for $x \in S^{ss}$.
\end{proof}

\end{document}
